% page 383 du fac-simile (image p-382 du scan)
% bloc 180.3 x 237.7 mm ; marge gauche 21.7 mm
\pgc{0.000mm}{0.000mm}{
\l{\cel{60}375}
\l{}
\l{\cel{5}\sou{minora}. I. (\sou{yurocienco}). Qua ne atingis ankore la evo lego-}
\l{\cel{8}stipulita koncerne la yuro agar libere pri sua persono e sua}
\l{\cel{8}havaji. - II. (\sou{muziko)}. Intervalo ek tri noti,qua kontenas}
\l{\cel{8}un tono ed un mi-tono, o de sis noti,qua kontenas tri toni}
\l{\cel{8}e du mi-toni. - III. (logiko). (en\cel{1}\sou{silogismo}). La subjekto}
\l{\cel{8}di la konkluzo (pro ke lu esas min ampla kam la atributo).-}
\l{\cel{8}DEFIS.}
\l{}
\l{\cel{5}\sou{minoritato}. La quanto minim granda de voci che asemblitaro\cel{1}\sou{qua}}
\l{\cel{8}esas votanta pri lo jus deliberita, DEFIS.}
\l{}
\l{\cel{5}minto.(\sou{bot.)}\cel{2}Planto aromatoza, ek la familio \textquotedbl{}labiacei\textquotedbl{}. L.}
\l{\cel{8}mentha. - DEFIRS.}
\l{}
\l{\cel{5}minucio. Detalo mikrega qua ne meritas ke onu haltez por lu (pro}
\l{\cel{8}lu). - DEFIS.}
\l{}
\l{\cel{5}minus. (arit., algeb.) Signo qua, avan quanto, signifikas ke olca}
\l{\cel{8}esas sustraktenda. DEIRS.}
\l{}
\l{\cel{5}minuskulo.(pri la literi di la alfabeto).Mikra litero (o tipo),}
\l{\cel{8}kompare a mayuskulo, e ye qua konsistas la maxim multa literi}
\l{\cel{8}di la vorti skribita. - DEFIS.}
\l{}
\l{\cel{5}minuto. I. Sisadekima parto di horo. - II. Sisadekima parto di un}
\l{\cel{8}grado di cirklo. - DEFIRS.}
\l{}
\l{\cel{5}mioceno. (geol.)\cel{2}Un ek la quar granda faki di la ero terciara}
\l{\cel{8}(olta qua sucedas la strati oligocena, e sur qua jacas la stra-}
\l{\cel{8}ti pliocena). - DEFIS.}
\l{}
\l{\cel{5}miologio. La anatomio di la muskuli. - DEFIS.}
\l{}
\l{\cel{5}miopa. (patol.) Di qua la vido-porteo esas kurta. - DEFIS.}
\l{}
\l{\cel{5}miozoto. (bot.) Planto di qua la floro esas ciel-blua, rozea,}
\l{\cel{8}e quan, en landi diversa diEuropa, onu nomizas poezioze}
\l{\cel{8}\textquotedbl{}Ne-obliviez-me\textquotedbl{}. - L.\cel{1}\sou{myosotis}. - EFISL.}
\l{}
\l{\cel{5}miro. (bot.)\cel{2}Gumo-rezino aromoza. - DEFIS.}
\l{}
\l{\cel{5}\sou{mirabelo}. Mikra pruno flava.}
\l{}
\l{\cel{5}\sou{mirajo}. Optik-iluziono - refrakto qua, en la landi di qui la}
\l{\cel{8}atmosfero esas varmega, igas vidar kom proxima la imaji di}
\l{\cel{8}objekti tre fora. - (\sou{metaf}.:Iluziono seduktoza). - EFIR}
\l{}
\l{\cel{5}\sou{miraklo}. Fakto supernatura, qua eventas violace di la +leyi}
\l{\cel{8}(legi naturala). - DEFIR.}
\l{}
\l{\cel{5}\sou{miriado.}\cel{2}10.000.}
\l{}
\l{\cel{5}\sou{miriametro}. 10.000 metri.}
}
